\chapter{Resultados e discussão}\label{chap:resultados}
Este capítulo é um roteiro de apresentação dos resultados. Em trabalhos reais, as figuras e tabelas devem refletir dados efetivamente obtidos e conter discussão compatível com os objetivos definidos.

\section{Organização dos experimentos}
Defina cenários, parâmetros de controle e critérios de comparação. Indique quais resultados são quantitativos, qualitativos ou exploratórios.

\section{Apresentação e análise}
A Tabela~\ref{tab:resultados} demonstra como organizar métricas, mas seus valores são fictícios. Em uma pesquisa real, substitua-a por dados rastreáveis, interprete tendências e discuta o significado físico e computacional dos resultados.

\section{Limitações}
Registre limites de generalização, hipóteses simplificadoras, dependência de parâmetros e eventuais diferenças entre configurações experimentais. Evite afirmar superioridade geral de um método a partir de um conjunto restrito de casos.
